\documentclass[twoside,10pt]{../jmlr} 
% \documentclass{colt2020} % Include author names

% The following packages will be automatically loaded:
% amsmath, amssymb, natbib, graphicx, url, algorithm2e

\allowdisplaybreaks
\newcommand{\jmlrBlackBox}{\rule{1.5ex}{1.5ex}}
\providecommand{\BlackBox}{\jmlrBlackBox}
\newcommand{\jmlrQED}{\hfill\jmlrBlackBox\par\bigskip}
\usepackage{graphicx}
\makeatletter
 \let\Ginclude@graphics\@org@Ginclude@graphics 
\makeatother

%\newcommand{\rb}[1]{\textcolor{magenta}{RB:#1}}

%\pdfoutput=1

%  \newcommand{\jmlrBlackBox}{\rule{1.5ex}{1.5ex}}
%    \providecommand{\BlackBox}{\jmlrBlackBox}
%  \newcommand{\jmlrQED}{\hfill\jmlrBlackBox\par\bigskip}


%\allowdisplaybreaks

\title{Notes on grifting by optimal transport}
\usepackage{times}
\usepackage[utf8]{inputenc} % allow utf-8 input
\usepackage[T1]{fontenc}    % use 8-bit T1 fonts
\usepackage{hyperref}       % hyperlinks
\usepackage{url}            % simple URL typesetting
\usepackage{booktabs}       % professional-quality tables
\usepackage{amsfonts}       % blackboard math symbols
\usepackage{nicefrac}       % compact symbols for 1/2, etc.
\usepackage{microtype}      % microtypography
\usepackage{xspace}
\usepackage{amsmath}


\usepackage{mathtools}  
\usepackage{amsmath}
\usepackage{amssymb}
\usepackage{tabulary}
\usepackage{booktabs}

\usepackage{bm}


\newtheorem{assumption}{Assumption}
\usepackage{algorithm}
\usepackage{algpseudocode}
\usepackage{algpascal}
\usepackage{comment}
\usepackage{selectp}




%\newcommand{\dirac}[2]{\delta_{#1}(#2)}
\newcommand{\kprimal}{K}
\newcommand{\dirac}[1]{\delta_{#1}}
\newcommand{\ustar}{u^\star}
\newcommand{\vstar}{v^\star}
\newcommand{\psistar}{\psi^\star}
\newcommand{\Vstar}{V^\star}
\newcommand{\Xstar}{X^\star}
\newcommand{\Ystar}{Y^\star}
\newcommand{\Vhat}{\wh{V}}
\newcommand{\pistar}{\pi^\star}
\newcommand{\pibar}{\overline{\pi}}
\newcommand{\mustar}{\mu^\star}
\newcommand{\hmu}{\wh{\mu}}

\newcommand{\Hplus}{H^+}
\newcommand{\Htest}{H_{\text{test}}}

\newcommand{\muout}{\mu_{\text{out}}}

\newcommand{\mupi}{\mu^\pi}
\newcommand{\nupi}{\nu^\pi}
\newcommand{\dpi}{d^\pi}
\newcommand{\hy}{\wh{y}}
\newcommand{\hx}{\wh{x}}

\newcommand{\piout}{\pi_{\text{out}}}

\newcommand{\tPi}{\wt{\Pi}}
\newcommand{\tpi}{\wt{\pi}}
\newcommand{\tmu}{\wt{\mu}}
\newcommand{\tc}{\wt{c}}
\newcommand{\bmu}{\overline{\mu}}
\newcommand{\bmuX}{\overline{\mu}_\X}
\newcommand{\HH}[2]{\mathcal{H}\pa{#1\middle\|#2}}
\newcommand{\HHX}[2]{\mathcal{H}_{\X}\pa{#1\middle\|#2}}

\newcommand{\bone}{\bm{1}}
\newcommand{\diag}{\mathrm{diag}}
\newcommand{\dd}{\mathrm{d}}
\newcommand{\nux}{\nu_\X}
\newcommand{\nuy}{\nu_\Y}

\newcommand{\bPi}{\overline{\Pi}}

\newcommand{\X}{\mathcal{X}}
\newcommand{\U}{\mathcal{U}}
\newcommand{\XI}{\X^{\infty}}
\newcommand{\PX}{\mathcal{P}_X}
\newcommand{\PY}{\mathcal{P}_Y}
\newcommand{\Y}{\mathcal{Y}}
\newcommand{\YI}{\mathcal{Y}^{\infty}}
\newcommand{\V}{\mathcal{V}}
\newcommand{\LL}{\mathcal{L}}
\newcommand{\GG}{\mathcal{G}}
\newcommand{\hGG}{\wh{\GG}}
\newcommand{\BB}{\mathcal{B}}
\newcommand{\TT}{\mathcal{T}}
\newcommand{\TTpi}{\mathcal{T}^{\pi}}
\newcommand{\TTpik}{\mathcal{T}^{\pi_k}}
\newcommand{\TTP}{\mathcal{T}_{P}}
\newcommand{\TTQ}{\mathcal{T}_{Q}}
\newcommand{\TTX}{\mathcal{T}_{\X}}
\newcommand{\TTY}{\mathcal{T}_{\Y}}
\newcommand{\BX}{\mathcal{B}_\mathcal{X}}
\newcommand{\BY}{\mathcal{B}_\mathcal{Y}}
\newcommand{\SX}{\mathcal{S}_\mathcal{X}}
\newcommand{\SY}{\mathcal{S}_\mathcal{Y}}
\newcommand{\PPX}{P_\mathcal{X}}
\newcommand{\PPY}{P_\mathcal{Y}}
\newcommand{\nuzerox}{\nu_{0,\X}}
\newcommand{\nuzeroy}{\nu_{0,\Y}}


\newcommand{\F}{\mathcal{F}}
\newcommand{\B}{\mathcal{B}}
\newcommand{\FX}{\mathcal{F}_{\mathcal{X}}}
\newcommand{\FXI}{\mathcal{F}_{\mathcal{X}}^{\infty}}
\newcommand{\FY}{\mathcal{F}_{\mathcal{Y}}}
\newcommand{\FYI}{\mathcal{F}_{\mathcal{Y}}^{\infty}}
% \newcommand{\A}{\mathcal{A}}
\newcommand{\N}{\mathcal{N}}
% \newcommand{\I}{\mathcal{I}}
\newcommand{\real}{\mathbb{R}}
\newcommand{\Dw}{\mathcal{D}}
\newcommand{\Sw}{\mathcal{S}}
\newcommand{\Pw}{\mathcal{P}}
\newcommand{\Rw}{\mathcal{R}}
\newcommand{\Ww}{\mathcal{W}}
\newcommand{\DD}[2]{\mathcal{D}\pa{#1\middle\|#2}}
\newcommand{\DDh}[2]{\mathcal{B}_{\bH}\pa{#1\middle\|#2}}
\newcommand{\DDp}[2]{\mathcal{D}_{\ell_p}\pa{#1\middle\|#2}}
\newcommand{\DDLC}[2]{\mathcal{D}_{\mathrm{LC}}\pa{#1\middle\|#2}}
\newcommand{\tDD}[2]{\wt{\mathcal{D}}\pa{#1\middle\|#2}}
\newcommand{\DDR}[2]{\mathcal{D}_R\pa{#1\middle\|#2}}
\newcommand{\DDsigma}[2]{\mathcal{D}_{\sigma}\pa{#1\middle\|#2}}
\newcommand{\DDKL}[2]{\mathcal{D}_{\mathrm{KL}}\pa{#1\middle\|#2}}
\newcommand{\DDxy}[2]{\mathcal{D}_{\mathrm{KL}}^{xy}\pa{#1\middle\|#2}}
\newcommand{\kldiv}[2]{d_{\mathrm{kl}}\pa{#1\middle\|#2}}
\newcommand{\DDPhi}[2]{\mathcal{B}_\Phi\pa{#1\middle\|#2}}
\newcommand{\DDalpha}[2]{\mathcal{D}_{\alpha}\pa{#1\middle\|#2}}
\newcommand{\DDpsi}[2]{\mathcal{D}_{\psi}\pa{#1\middle\|#2}}
\newcommand{\DDchi}[2]{\mathcal{D}_{\chi^2}\pa{#1\middle\|#2}}
\newcommand{\DDphi}[2]{\mathcal{D}_{\varphi}\pa{#1\middle\|#2}}
\newcommand{\DDPhiS}[2]{\mathcal{D}_{\Phi_S}\pa{#1\middle\|#2}}
\newcommand{\DDtPhi}[2]{\mathcal{D}_{\tPhi}\pa{#1\middle\|#2}}
\newcommand{\DDS}[2]{\mathcal{D}^S\pa{#1\middle\|#2}}
\newcommand{\DDc}[3]{\mathcal{D}^{#3}\pa{#1\middle\|#2}}
\newcommand{\tDDS}[2]{\wt{\mathcal{D}}^S\pa{#1\middle\|#2}}
\newcommand{\DDf}[2]{\mathcal{D}_\varphi\pa{#1\middle\|#2}}
\newcommand{\OO}{\mathcal{O}}
\newcommand{\tOO}{\wt{\OO}}
\newcommand{\trace}[1]{\mbox{trace}\left(#1\right)}
\newcommand{\II}[1]{\mathbb{I}_{\left\{#1\right\}}}
\newcommand{\PP}[1]{\mathbb{P}\left[#1\right]}
\newcommand{\PPZ}[1]{\mathbb{P}_Z\left[#1\right]}
\newcommand{\EE}[1]{\mathbb{E}\left[#1\right]}
\newcommand{\EEP}[1]{\mathbb{E}_P\left[#1\right]}
\newcommand{\EES}[1]{\mathbb{E}_S\left[#1\right]}
\newcommand{\EESb}[1]{\mathbb{E}_S\bigl[#1\bigr]}
\newcommand{\EESB}[1]{\mathbb{E}_S\Bigl[#1\Bigr]}
\newcommand{\EEZ}[1]{\mathbb{E}_Z\left[#1\right]}
\newcommand{\EEb}[1]{\mathbb{E}\bigl[#1\bigr]}
\newcommand{\EEtb}[1]{\mathbb{E}_t\bigl[#1\bigr]}
\newcommand{\EXP}{\mathbb{E}}
\newcommand{\EEpi}[1]{\mathbb{E}_{\pi}\left[#1\right]}
\newcommand{\PPpi}[1]{\mathbb{P}_{\pi}\left[#1\right]}
\newcommand{\EEs}[2]{\mathbb{E}_{#2}\left[#1\right]}
\newcommand{\EEu}[1]{\mathbb{E}_{\u}\left[#1\right]}
\newcommand{\EEst}[1]{\mathbb{E}_{*}\left[#1\right]}
\newcommand{\EEo}[1]{\mathbb{E}_{0}\left[#1\right]}
\newcommand{\PPt}[1]{\mathbb{P}_t\left[#1\right]}
\newcommand{\EEt}[1]{\mathbb{E}_t\left[#1\right]}
\newcommand{\PPi}[1]{\mathbb{P}_i\left[#1\right]}
\newcommand{\EEi}[1]{\mathbb{E}_i\left[#1\right]}
\newcommand{\PPc}[2]{\mathbb{P}\left[#1\left|#2\right.\right]}
\newcommand{\PPs}[2]{\mathbb{P}_{#2}\left[#1\right]}
\newcommand{\PPcs}[3]{\mathbb{P}_{#3}\left[#1\left|#2\right.\right]}
\newcommand{\PPct}[2]{\mathbb{P}_t\left[#1\left|#2\right.\right]}
\newcommand{\PPcc}[2]{\mathbb{P}\left[\left.#1\right|#2\right]}
\newcommand{\PPcct}[2]{\mathbb{P}_t\left[\left.#1\right|#2\right]}
\newcommand{\PPcci}[2]{\mathbb{P}_i\left[\left.#1\right|#2\right]}
\newcommand{\EEc}[2]{\mathbb{E}\left[#1\left|#2\right.\right]}
\newcommand{\EEcc}[2]{\mathbb{E}\left[\left.#1\right|#2\right]}
\newcommand{\EEcs}[3]{\mathbb{E}_{#3}\left[\left.#1\right|#2\right]}
\newcommand{\EEcct}[2]{\mathbb{E}_t\left[\left.#1\right|#2\right]}
\newcommand{\EEcci}[2]{\mathbb{E}_i\left[\left.#1\right|#2\right]}
%\renewcommand{\th}{\ensuremath{^{\mathrm{th}}}}
\def\argmin{\mathop{\mbox{ arg\,min}}}
\def\argmax{\mathop{\mbox{ arg\,max}}}
\newcommand{\ra}{\rightarrow}

\newcommand{\siprod}[2]{\langle#1,#2\rangle}
\newcommand{\iprod}[2]{\left\langle#1,#2\right\rangle}
\newcommand{\Siprod}[2]{\left\langle#1,#2\right\rangle_{|S}}
\newcommand{\biprod}[2]{\bigl\langle#1,#2\bigr\rangle}
\newcommand{\Biprod}[2]{\Bigl\langle#1,#2\Bigr\rangle}
\newcommand{\norm}[1]{\left\|#1\right\|}
\newcommand{\bnorm}[1]{\bigl\|#1\bigr\|}
\newcommand{\onenorm}[1]{\norm{#1}_1}
\newcommand{\twonorm}[1]{\norm{#1}_2}
\newcommand{\infnorm}[1]{\norm{#1}_\infty}
\newcommand{\tvnorm}[1]{\norm{#1}_{\mathrm{TV}}}
\newcommand{\Hnorm}[1]{\norm{#1}_{\mathcal{H}}}
\newcommand{\ev}[1]{\left\{#1\right\}}
% \newcommand{\abs}[1]{\left|#1\right|}
\newcommand{\babs}[1]{\bigl|#1\bigr|}
\newcommand{\pa}[1]{\left(#1\right)}
\newcommand{\bpa}[1]{\bigl(#1\bigr)}
\newcommand{\Bpa}[1]{\Bigl(#1\Bigr)}
\newcommand{\BPA}[1]{\Biggl(#1\Biggr)}
\newcommand{\sign}{\mbox{sign}}

\newcommand{\wh}{\widehat}
\newcommand{\wt}{\widetilde}


\newcommand{\hr}{\wh{r}}
\newcommand{\hf}{\wh{f}}
\newcommand{\hs}{\wh{s}}
\newcommand{\bA}{\overline{A}}
\newcommand{\deff}{d_{\text{eff}}}

\newcommand{\Cinf}{C_\infty}
\newcommand{\tfs}{\wt{f}_{\sigma}}
\newcommand{\tQ}{\wt{Q}}
\newcommand{\tW}{\wt{W}}
\newcommand{\tR}{\wt{R}}
\newcommand{\tPhi}{\wt{\Phi}}
\newcommand{\tZ}{\wt{Z}}
\newcommand{\tS}{\wt{S}}
\newcommand{\tD}{\wt{D}}
\newcommand{\tw}{\wt{w}}
\newcommand{\bw}{\overline{w}}
\newcommand{\hX}{\wh{X}}
\newcommand{\hS}{\wh{\Sigma}}
\newcommand{\hw}{\wh{w}}
\newcommand{\bX}{\overline{X}}
\newcommand{\bY}{\overline{Y}}
\newcommand{\bS}{\overline{S}}
\newcommand{\bZ}{\overline{Z}}
\newcommand{\bW}{\overline{W}}
\newcommand{\bP}{\overline{P}}
\newcommand{\bDelta}{\overline{\Delta}}
\newcommand{\Pj}{P_X\otimes P_Y}
\newcommand{\barf}{\overline{f}}
\newcommand{\hQ}{\wh{Q}}
\newcommand{\hP}{\wh{P}}
\newcommand{\tP}{\wt{P}}
\newcommand{\transpose}{^\mathsf{\scriptscriptstyle T}}

\newcommand{\bloss}{\overline{\ell}}
\newcommand{\hloss}{\wh{\ell}}
\newcommand{\loss}{\ell}
\newcommand{\gen}{\textup{gen}}

\usepackage{todonotes}
\definecolor{PalePurp}{rgb}{0.66,0.57,0.66}
\newcommand{\todoG}[1]{\todo[color=PalePurp!30]{\textbf{Grg:} #1}}
\newcommand{\todoGI}[1]{\todo[inline,color=PalePurp!30]{#1}}
% % \newcommand{\reddd}[1]{}
\newcommand{\redd}[1]{\textcolor{red}{#1}}
\newcommand{\grg}[1]{\textcolor{red}{[\textbf{Grg:} #1]}}
\newcommand{\lud}[1]{\textcolor{red}{[\textbf{Ludovic:} #1]}}
\newcommand{\ser}[1]{\textcolor{red}{[\textbf{Sergio:} #1]}}
\newcommand{\gbr}[1]{\textcolor{orange}{[\textbf{Gbr:} #1]}}

\newcommand{\regret}{\mathrm{regret}}

% \newcommand{\qed}{\hfill\BlackBox\\[2mm]}


\newcommand{\hL}{\wh{L}}
\newcommand{\snr}{\mathtt{SNR}}

\newcommand{\alg}{\mathcal{A}}
\newcommand{\Zw}{\mathcal{Z}}
\newcommand{\bL}{\overline{L}}
%\newcommand{\bH}{\overline{H}}
\newcommand{\bH}{h}
\renewcommand{\th}{\wt{h}}
\newcommand{\Law}{\text{Law}}

\usepackage[normalem]{ulem}
\usepackage{enumitem}

% Put your macros here

\newcommand\am[1]{{\color{red}(\textbf{AM:} #1)}}

\newcommand\pSrc{p_{\text{src}}}
\newcommand\Pm{\mathbb{P}}
\newcommand\supp{\text{supp}}
\newcommand\opt{\text{OPT}}

\newcommand\src{\nu_{\text{src}}}
\newcommand\tgt{\nu_{\text{tgt}}}
\newcommand\dSrc{\rho_{\text{src}}}
\newcommand\dTgt{\rho_{\text{tgt}}}
\newcommand\Wc{\mathcal{W}}
\newcommand\fwd{\mathcal{F}}
\newcommand\bwd{\mathcal{B}}
\newcommand\Sc{\mathcal{S}}
% \newcommand\norm[1]{\left\|#1\right\|}
\newcommand\Lc{\mathcal{L}}
\newcommand\hLc{\wh{\Lc}}
%\newcommand\Cc{\mathcal{C}}

% \usepackage{etoc}
% \usepackage{authblk}
% \usepackage{microtype}
% \usepackage{graphicx}
% %\usepackage{subcaption}
% \usepackage[T1]{fontenc}
% \usepackage[english]{babel}
% \usepackage[utf8]{inputenc}
% \usepackage{booktabs}
% \usepackage{array}
% \usepackage{amsmath}
% \usepackage{amssymb}
% \usepackage{mathtools}
% \usepackage{amsthm}
% \usepackage{bm}
% \usepackage{thmtools}
% \usepackage{thm-restate}
% \usepackage{hyperref}
% \hypersetup{
%     colorlinks = true,
%     linkcolor = blue,
%     citecolor=blue,
%     urlcolor=black,
% }
% \usepackage{xcolor}
% \definecolor{mblue}{rgb}{0.0,0.0,1.0}
% \definecolor{mgreen}{rgb}{0.0,0.501960784314,0.0}
% \definecolor{mred}{rgb}{1.0,0.0,0.0}
% \usepackage[capitalize]{cleveref}
% \usepackage[ruled]{algorithm2e}
% \usepackage[noend]{algpseudocode}
% \renewcommand{\algorithmiccomment}[1]{\hfill$\triangleright~${\color{gray}#1}}
% \usepackage{bbm}
% \usepackage{xcolor}
% \usepackage{parskip}
% \usepackage{setspace}
% %\usepackage{etoolbox}
% \usepackage{comment}
% \usepackage{caption}
% \usepackage{enumitem}
% \usepackage{wrapfig}
% 
% \newtheorem{theorem}{Theorem}[section]
% \newtheorem{example}{Example}[section]
% \newtheorem{remark}{Remark}[section]
% \newtheorem{corollary}{Corollary}[section]
% \newtheorem{definition}{Definition}[section]
% \newtheorem{lemma}{Lemma}[section]
% \newtheorem{claim}{Claim}[section]
% \newtheorem{observation}{Observation}[section]
% \newtheorem{proposition}{Proposition}[section]
% \newtheorem{assumption}{Assumption}[section]
% \newtheorem{question}{Question}
% \newtheorem{itheorem}{Informal Theorem}[section]

\newcommand\R{\mathbb{R}}
\newcommand\inner[2]{\left\langle #1, #2 \right\rangle}
\newcommand\indicator[1]{\mathbbm{1}\left\{ #1 \right\}}
\newcommand\Simplex[1]{\Delta_{#1}}

\newcommand\Ac{\mathcal{A}}
\newcommand\Bc{\mathcal{B}}
\newcommand\Cc{\mathcal{C}}
\newcommand\Dc{\mathcal{D}}
\newcommand\Fc{\mathcal{F}}
\newcommand\Xc{\mathcal{X}}
\newcommand\Yc{\mathcal{Y}}
\newcommand\Mc{\mathcal{M}}
\newcommand\Tc{\mathcal{T}}
%\newcommand\E{\mathbb{E}}
\DeclareMathOperator*{\E}{\mathbb{E}}
\newcommand\Dkl{D_{\text{KL}}}
\newcommand\pre{p^{\text{pre}}}
\newcommand\dHat{\hat{d}}
\newcommand\sHat{\hat{s}}
\newcommand\pHat{\hat{p}}

\newcommand\OTilde{\tilde{O}}
\newcommand\biggVert{~\bigg\vert~}
\newcommand\bigVert{~\big\vert~}
\newcommand\normVert{~\vert~}
\newcommand\rHat{\hat{r}}
\newcommand\floor[1]{\lfloor #1 \rfloor}
\newcommand\ceil[1]{\lceil #1 \rceil}
\newcommand{\bc}[1]{\left\{{#1}\right\}}
\newcommand{\br}[1]{\left({#1}\right)}
\newcommand{\bs}[1]{\left[{#1}\right]}
\newcommand{\abs}[1]{\left\lvert #1 \right\rvert}
\newcounter{protocol}
\makeatletter
\newenvironment{protocol}[1][htb]{%
  \let\c@algorithm\c@protocol
  \renewcommand{\ALG@name}{Protocol}
  \begin{algorithm}[#1]%
  }{\end{algorithm}
}

\makeatletter
\providecommand*{\bigcupdot}{%
  \mathop{%
    \vphantom{\bigcup}%
    \mathpalette\@bigcupdot{}%
  }%
}
\newcommand*{\@bigcupdot}[2]{%
  \ooalign{%
    $\m@th#1\bigcup$\cr
    \sbox0{$#1\bigcup$}%
    \dimen@=\ht0 %
    \advance\dimen@ by -\dp0 %
    \sbox0{\scalebox{1.5}{$\m@th#1\cdot$}}%
    \advance\dimen@ by -\ht0 %
    \dimen@=.5\dimen@
    \hidewidth\raise\dimen@\box0\hidewidth
  }%
}
\makeatother



% Operators

\def\x{\times}

\newcommand{\toas}{\xrightarrow{\text{a.s.}}}
\newcommand{\todis}{\xrightarrow{\text{d}}}
%\newcommand{\supp}{\text{supp}}
\newcommand{\doubint}{\!\int\!\!\!\int}

%\DeclareMathOperator*{\argmax}{argmax}
%\DeclareMathOperator*{\argmin}{argmin}
\DeclareMathOperator{\Trace}{Tr}
\DeclareMathOperator{\Tr}{Tr}
%\DeclareMathOperator{\trace}{tr}
\DeclareMathOperator{\support}{supp}
\DeclareMathOperator{\osc}{osc}
\DeclareMathOperator{\Sp}{sp}
\DeclareMathOperator{\sgn}{sgn}
\DeclareMathOperator{\interior}{int}
\DeclareMathOperator{\Id}{Id}
\DeclareMathOperator{\idmat}{\boldmath{I}}
%\DeclareMathOperator{\diag}{diag}
\DeclareMathOperator{\dimension}{dim}
\DeclareMathOperator{\diam}{diam}


\def\mybigtimes{\mathop{\mathchoice{%display:
   \vcenter{\hbox to10bp{\vrule height15bp width0pt \pdfliteral{
   q 1 J .8 w 0 1 m 10 14 l S 0 14 m 10 1 l S Q
}\hss}}}{%text:
   \vcenter{\hbox to10bp{\kern1bp\vrule height10bp width0pt \pdfliteral{
   q 1 J .65 w 0 0 m 8 10 l S 0 10 m 8 0 l S Q
}\hss}}}{\times}{\times}%script, scriptscript not defined
}}

\DeclareMathOperator*{\bigtimes}{\mybigtimes}

% Delimiters

%\DeclarePairedDelimiter\ceil{\lceil}{\rceil}
%\DeclarePairedDelimiter\floor{\lfloor}{\rfloor}

%\newcommand{\abs}[1]{\left\lvert #1 \right\rvert}
\newcommand{\sabs}[1]{\lvert #1\rvert}

\newcommand{\card}[1]{\abs{#1}}
\newcommand{\scard}[1]{\sabs{#1}}

%\newcommand{\norm}[1]{\left\lVert#1\right\rVert}
\newcommand{\snorm}[1]{\lVert#1\rVert}



% Text

\def\iid{\overset{\mathrm{iid}}{\sim}}


\makeatletter
\newcommand{\vvast}{\bBigg@{3}}
\newcommand{\vast}{\bBigg@{4}}
\newcommand{\Vast}{\bBigg@{5}}
\makeatother



% Characters

\newcommand{\NN}{\mathbb{N}}
\newcommand{\ZZ}{\mathbb{Z}}
\newcommand{\RR}{\mathbb{R}}
% \newcommand{\PP}{\mathbb{P}}
% \newcommand{\EE}{\mathbb{E}}
% \newcommand{\II}{\mathbb{I}}
\newcommand{\FF}{\mathbb{F}}


\newcommand{\1}{\mathds{1}}


\def\ar{\mathrm{a}}
\def\ccr{\mathrm{c}} %special \cr is already used
\def\Rr{\mathrm{R}}
\def\Wr{\mathrm{W}}

\newcommand{\op}{{\rm op}}
\newcommand{\de}{\mathrm{d}}
\newcommand{\De}{\mathrm{D}}

\newcommand{\ve}{\varepsilon}
\newcommand{\eps}{\varepsilon}
\newcommand{\vp}{\varphi}

\newcommand{\tensor}{\otimes}


%%%%%%%%%%%%%%%%%%%%%%%%%%
% Alphabet
%%%%%%%%%%%%%%%%%%%%%%%%%%


%%% Blackboard Bold
    \def\Ab{\mathbb{A}}
    \def\Bb{\mathbb{B}}
    \def\Cb{\mathbb{C}}
    \def\Db{\mathbb{D}}
    \def\Eb{\mathbb{E}}
    \def\Fb{\mathbb{F}}
    \def\Hb{\mathbb{H}}
    \def\Gb{\mathbb{G}}
    \def\Ib{\mathbb{I}}
    \def\Jb{\mathbb{J}}
    \def\Lb{\mathbb{L}}
    \def\Kb{\mathbb{K}}
    \def\Mb{\mathbb{M}}
    \def\Nb{\mathbb{N}}
    \def\Ob{\mathbb{O}}
    \def\Pb{\mathbb{P}}
    \def\Qb{\mathbb{Q}}
    \def\Rb{\mathbb{R}}
    \def\Sb{\mathbb{S}}
    \def\Tb{\mathbb{T}}
    \def\Ub{\mathbb{U}}
    \def\Vb{\mathbb{V}}
    \def\Wb{\mathbb{W}}
    \def\Xb{\mathbb{X}}
    \def\Yb{\mathbb{Y}}
    \def\Zb{\mathbb{Z}}

%%% Calygraphic

    \def\Ac{\mathcal{A}}
    \def\Bc{\mathcal{B}}
    \def\Cc{\mathcal{C}}
    \def\Dc{\mathcal{D}}
    \def\Ec{\mathcal{E}}
    \def\Fc{\mathcal{F}}
    \def\Hc{\mathcal{H}}
    \def\Gc{\mathcal{G}}
    \def\Ic{\mathcal{I}}
    \def\Jc{\mathcal{J}}
    \def\Lc{\mathcal{L}}
    \def\Kc{\mathcal{K}}
    \def\Mc{\mathcal{M}}
    \def\Nc{\mathcal{N}}
    \def\Oc{\mathcal{O}}
    \def\Pc{\mathcal{P}}
    \def\Qc{\mathcal{Q}}
    \def\Rc{\mathcal{R}}
    \def\Sc{\mathcal{S}}
    \def\Tc{\mathcal{T}}
    \def\Uc{\mathcal{U}}
    \def\Vc{\mathcal{V}}
    \def\Wc{\mathcal{W}}
    \def\Xc{\mathcal{X}}
    \def\Yc{\mathcal{Y}}
    \def\Zc{\mathcal{Z}}

% %%% Roman

    \def\Ar{\mathrm{A}}
    \def\Br{\mathrm{B}}
    \def\Cr{\mathrm{C}}
    \def\Dr{\mathrm{D}}
    \def\Er{\mathrm{E}}
    \def\Fr{\mathrm{F}}
    \def\Hr{\mathrm{H}}
    \def\Gr{\mathrm{G}}
    \def\Ir{\mathrm{I}}
    \def\Jr{\mathrm{J}}
    \def\Lr{\mathrm{L}}
    \def\Kr{\mathrm{K}}
    \def\Mr{\mathrm{M}}
    \def\Nr{\mathrm{N}}
    \def\Or{\mathrm{O}}
    \def\Pr{\mathrm{P}}
    \def\Qr{\mathrm{Q}}
    \def\Rr{\mathrm{R}}
    \def\Sr{\mathrm{S}}
    \def\TTr{\mathrm{T}} % trace
    \def\Ur{\mathrm{U}}
    \def\Vr{\mathrm{V}}
    \def\Wr{\mathrm{W}}
    \def\Xr{\mathrm{X}}
    \def\Yr{\mathrm{Y}}
    \def\Zr{\mathrm{Z}}

    \def\ar{\mathrm{a}}
    \def\br{\mathrm{b}}
    \def\ccr{\mathrm{c}} % \cr already used by some internals
    \def\dr{\mathrm{d}}
    \def\er{\mathrm{e}}
    \def\fr{\mathrm{f}}
    \def\hr{\mathrm{g}}
    \def\gr{\mathrm{h}}
    \def\ir{\mathrm{i}}
    \def\jr{\mathrm{j}}
    \def\kr{\mathrm{k}}
    \def\lr{\mathrm{l}}
    \def\mr{\mathrm{m}}
    \def\nr{\mathrm{n}}
    \def\oor{\mathrm{o}} %\or logic
    \def\pr{\mathrm{p}}
    \def\qr{\mathrm{q}}
    \def\rr{\mathrm{r}}
    \def\sr{\mathrm{s}}
    \def\tr{\mathrm{t}}
    \def\ur{\mathrm{u}}
    \def\vr{\mathrm{v}}
    \def\wr{\mathrm{w}}
    \def\xr{\mathrm{x}}
    \def\yr{\mathrm{y}}
    \def\zr{\mathrm{z}}

%%% Script

    \def\As{\mathscr{A}}
    \def\Bs{\mathscr{B}}
    \def\Cs{\mathscr{C}}
    \def\Ds{\mathscr{D}}
    \def\Es{\mathscr{E}}
    \def\Fs{\mathscr{F}}
    \def\Hs{\mathscr{H}}
    \def\Gs{\mathscr{G}}
    \def\Is{\mathscr{I}}
    \def\Js{\mathscr{J}}
    \def\Ls{\mathscr{L}}
    \def\Ks{\mathscr{K}}
    \def\Ms{\mathscr{M}}
    \def\Ns{\mathscr{N}}
    \def\Os{\mathscr{O}}
    \def\Ps{\mathscr{P}}
    \def\Qs{\mathscr{Q}}
    \def\Rs{\mathscr{R}}
    \def\Ss{\mathscr{S}}
    \def\Ts{\mathscr{T}}
    \def\Us{\mathscr{U}}
    \def\Vs{\mathscr{V}}
    \def\Ws{\mathscr{W}}
    \def\Xs{\mathscr{X}}
    \def\Ys{\mathscr{Y}}
    \def\Zs{\mathscr{Z}}

    \def\V{\Vs}

%%% Fraktur

\def\Af{\mathfrak{A}}
    \def\Bf{\mathfrak{B}}
    \def\Cf{\mathfrak{C}}
    \def\Df{\mathfrak{D}}
    \def\Ef{\mathfrak{E}}
    \def\Ff{\mathfrak{F}}
    \def\Hf{\mathfrak{H}}
    \def\Gf{\mathfrak{G}}
    %\def\If{\mathfrak{I}}
    \def\Jf{\mathfrak{J}}
    \def\Lf{\mathfrak{L}}
    \def\Kf{\mathfrak{K}}
    \def\Mf{\mathfrak{M}}
    \def\Nf{\mathfrak{N}}
    \def\Of{\mathfrak{O}}
    \def\Pf{\mathfrak{P}}
    \def\Qf{\mathfrak{Q}}
    \def\Rf{\mathfrak{R}}
    \def\Sf{\mathfrak{S}}
    \def\Tf{\mathfrak{T}}
    \def\Uf{\mathfrak{U}}
    \def\Vf{\mathfrak{V}}
    \def\Wf{\mathfrak{W}}
    \def\Xf{\mathfrak{X}}
    \def\Yf{\mathfrak{Y}}
    \def\Zf{\mathfrak{Z}}


    \def\af{\mathfrak{a}}
    \def\bbf{\mathfrak{b}} % \bf boldface
    \def\cf{\mathfrak{c}}
    \def\df{\mathfrak{d}}
    \def\ef{\mathfrak{e}}
    \def\ff{\mathfrak{f}}
    \def\hf{\mathfrak{g}}
    \def\gf{\mathfrak{h}}
    \def\iif{\mathfrak{i}} % \if conditional statements
    \def\jf{\mathfrak{j}}
    \def\kf{\mathfrak{k}}
    \def\lf{\mathfrak{l}}
    \def\mf{\mathfrak{m}}
    \def\nf{\mathfrak{n}}
    \def\of{\mathfrak{o}}
    \def\pf{\mathfrak{p}}
    \def\qf{\mathfrak{q}}
    \def\rf{\mathfrak{r}}
    \def\sf{\mathfrak{s}}
    \def\tf{\mathfrak{t}}
    \def\uf{\mathfrak{u}}
    \def\vf{\mathfrak{v}}
    \def\wf{\mathfrak{w}}
    \def\xf{\mathfrak{x}}
    \def\yf{\mathfrak{y}}
    \def\zf{\mathfrak{z}}

% %%% Bold Face

    \def\Abf{\mathbf{A}}
    \def\Bbf{\mathbf{B}}
    \def\Cbf{\mathbf{C}}
    \def\Dbf{\mathbf{D}}
    \def\Ebf{\mathbf{E}}
    \def\Fbf{\mathbf{F}}
    \def\Hbf{\mathbf{H}}
    \def\Gbf{\mathbf{G}}
    \def\Ibf{\mathbf{I}}
    \def\Jbf{\mathbf{J}}
    \def\Lbf{\mathbf{L}}
    \def\Kbf{\mathbf{K}}
    \def\Mbf{\mathbf{M}}
    \def\Nbf{\mathbf{N}}
    \def\Obf{\mathbf{O}}
    \def\Pbf{\mathbf{P}}
    \def\Qbf{\mathbf{Q}}
    \def\Rbf{\mathbf{R}}
    \def\Sbf{\mathbf{S}}
    \def\Tbf{\mathbf{T}}
    \def\Ubf{\mathbf{U}}
    \def\Vbf{\mathbf{V}}
    \def\Wbf{\mathbf{W}}
    \def\Xbf{\mathbf{X}}
    \def\Ybf{\mathbf{Y}}
    \def\Zbf{\mathbf{Z}}


    \def\abf{\mathbf{a}}
    \def\bbf{\mathbf{b}}
    \def\cbf{\mathbf{c}}
    \def\dbf{\mathbf{d}}
    \def\ebf{\mathbf{e}}
    \def\fbf{\mathbf{f}}
    \def\hbf{\mathbf{g}}
    \def\gbf{\mathbf{h}}
    \def\ibf{\mathbf{i}}
    \def\jbf{\mathbf{j}}
    \def\kbf{\mathbf{k}}
    \def\lbf{\mathbf{l}}
    \def\mbf{\mathbf{m}}
    \def\nbf{\mathbf{n}}
    \def\obf{\mathbf{o}}
    \def\pbf{\mathbf{p}}
    \def\qbf{\mathbf{q}}
    \def\rbf{\mathbf{r}}
    \def\sbf{\mathbf{s}}
    \def\tbf{\mathbf{t}}
    \def\ubf{\mathbf{u}}
    \def\vbf{\mathbf{v}}
    \def\wbf{\mathbf{w}}
    \def\xbf{\mathbf{x}}
    \def\ybf{\mathbf{y}}
    \def\zbf{\mathbf{z}}








%\PassOptionsToPackage{numbers, sort&compress}{natbib}

%\bibliographystyle{plainnat}
\setcitestyle{authoryear,open={(},close={)}}
\usepackage{algorithm2e}
%\usepackage{hyperref}       % hyperlinks
\usepackage{url}            % simple URL typesetting
\usepackage{booktabs}       % professional-quality tables
\usepackage{nicefrac}       % compact symbols for 1/2, etc.
\usepackage{microtype}      % microtypography
%\usepackage{xcolor}         % colors
%\usepackage{caption}        % captions
%\usepackage{subcaption}     % subfigures
\usepackage{multirow}       % multirow in tables
\usepackage{listings}       % code
\usepackage[english]{babel}




%\usepackage[dvipsnames]{xcolor}

% PGF
\usepackage{pgfplots}
\pgfplotsset{width=8cm,compat=1.15}

% Attempt to make hyperref and algorithmic work together better:
%\newcommand{\theHalgorithm}{\arabic{algorithm}}

%symbols 
\usepackage{amsmath,amssymb,amsfonts,amsopn,mathrsfs}
\usepackage{bm,dsfont}
\usepackage{mathtools}
\usepackage{bigints}
\usepackage{latexsym}
\usepackage{stmaryrd}       %weird symbols like square \cap and double brackets


%theorems
%\usepackage{thmtools}
%\usepackage{thm-restate}

% Cleveref..
\usepackage[capitalize,noabbrev]{cleveref}

% Explicit names for cleveref (fix environments being labeled as "Theorem")
\crefname{theorem}{Theorem}{Theorems}
\Crefname{theorem}{Theorem}{Theorems}
\crefname{lemma}{Lemma}{Lemmas}
\Crefname{lemma}{Lemma}{Lemmas}
\crefname{proposition}{Proposition}{Propositions}
\Crefname{proposition}{Proposition}{Propositions}
\crefname{corollary}{Corollary}{Corollaries}
\Crefname{corollary}{Corollary}{Corollaries}
\crefname{definition}{Definition}{Definitions}
\Crefname{definition}{Definition}{Definitions}
\crefname{remark}{Remark}{Remarks}
\Crefname{remark}{Remark}{Remarks}
\crefname{example}{Example}{Examples}
\Crefname{example}{Example}{Examples}
\crefname{assumption}{Assumption}{Assumptions}
\Crefname{assumption}{Assumption}{Assumptions}
\crefname{conjecture}{Conjecture}{Conjectures}
\Crefname{conjecture}{Conjecture}{Conjectures}
\crefname{axiom}{Axiom}{Axioms}
\Crefname{axiom}{Axiom}{Axioms}
\crefname{claim}{Claim}{Claims}
\Crefname{claim}{Claim}{Claims}
\crefname{observation}{Observation}{Observations}
\Crefname{observation}{Observation}{Observations}

% jmlr defines several theorem-like environments by reusing the 'theorem'
% counter. Wrap those existing environments so cleveref records the original
% environment type instead of collapsing everything to 'theorem'.
\makeatletter
\let\jmlr@orig@thm\@thm
\let\jmlr@cref@thmtype\@empty
\def\@thm#1#2{%
	\ifx\jmlr@cref@thmtype\@empty
		\refstepcounter{#1}%
	\else
		\refstepcounter[\jmlr@cref@thmtype]{#1}%
		\global\let\jmlr@cref@thmtype\@empty
	\fi
	\@ifnextchar[{\@ythm{#1}{#2}}{\@xthm{#1}{#2}}%
}
\newcommand{\jmlr@patchsharedtheorem}[1]{%
	\expandafter\let\csname jmlr@orig@env@#1\expandafter\endcsname\csname #1\endcsname
	\expandafter\def\csname #1\endcsname{%
		\gdef\jmlr@cref@thmtype{#1}%
		\csname jmlr@orig@env@#1\endcsname
	}%
}
\jmlr@patchsharedtheorem{lemma}
\jmlr@patchsharedtheorem{proposition}
\jmlr@patchsharedtheorem{corollary}
\jmlr@patchsharedtheorem{definition}
\jmlr@patchsharedtheorem{remark}
\jmlr@patchsharedtheorem{example}
\jmlr@patchsharedtheorem{conjecture}
\jmlr@patchsharedtheorem{axiom}
\makeatother
%%%%%%%%%%%%%%%%%%%%%%%%%%%%%%%%
% THEOREMS
%%%%%%%%%%%%%%%%%%%%%%%%%%%%%%%%
% \theoremstyle{plain}
% \newtheorem{theorem}{Theorem}[section]
% \newtheorem{proposition}[theorem]{Proposition}
% \newtheorem{lemma}[theorem]{Lemma}
% \newtheorem{corollary}[theorem]{Corollary}

% \theoremstyle{definition}
% %\newtheorem{definition}{Definition}
% \newtheorem{assumption}{Assumption}

% \crefname{assumption}{assumption}{assumptions}
% \theoremstyle{remark}
% \newtheorem{remark}{Remark}[section]
% \newtheorem{example}{Example}[section]

\usepackage{textcomp}


%General macros
\def\regret{\Rs}
\def\fourier{\mathsf{F}}
\def\reflection{\mathsf{R}}
\def\measures{\mathscr{M}}
\def\fset{\Theta}
\def\ffset{\Fs}

\def\confset{\Cc}
\def\entropy{\mathscr{H}}
\def\event{\Ec}

\def\LSC{\mathsf{LSC}}
\def\ubnorm{C}
\def\width{\beta}
\def\design{V}
\def\designl{\design^\lambda}
\def\entf{\Psi_{\mu,\nu}^\ve}

\def\Lip{\mathrm{Lip}}
% Notation

\def\statespace{\mathcal{S}}
\def\actionspace{\mathcal{A}}
\def\stateactionspace{\mathcal{X}}

\newcommand{\fop}{\Fc}
\newcommand{\bop}{\Bc}

\newcommand{\pushforward}{\vert_\sharp}
\newcommand{\primalvalue}{J_p}
\newcommand{\dualvalue}{J_d}
\newcommand{\pspaceomega}{\Omega}
\newcommand{\pspacef}{\Sigma}

\DeclareMathOperator{\bigcomp}{\bigcirc}

%project macros 






\author{
 \Name[{Gergely~Neu}]{Gergely Neu} \Email{gergely.neu@gmail.com}\\
 \addr Universitat Pompeu Fabra, Barcelona, Spain
 \AND
 \Name[{Lorenzo~Croissant}]{Lorenzo Croissant} \Email{lorenzo.croissant@upf.edu}\\
 \addr Universitat Pompeu Fabra, Barcelona, Spain
 }
 
\begin{document}
\maketitle

% \begin{abstract}
 
% \end{abstract}

\begin{keywords}%
  optimal transport, generative modeling, Markov chains
\end{keywords}


Consider an affine measurable space $\statespace\subset\mathbb{R}^d$. Fixing $(\pspaceomega,\pspacef,\Pb)$, $\pspaceomega=\statespace$, $\pspacef=\Bs(\statespace)$ the Borell $\sigma$-algebra of $\statespace$, and $\Pb$ a probability measure on $(\pspaceomega,\pspacef)$. Let $v:=(v_t)_{t\in\Nb}$ be a collection of $\pspacef$-measurable automorphisms on $\statespace$. For any random variable $X:\pspaceomega\to\statespace$, let $\Lc(X)$ denote the law of $X$ under $\Pb$.

The dynamical system of a drifting model \citep{deng2026generative} evolves the prediction $X$ during training according to the following Markov chain:
\begin{align}
    X_{t+1}^\velocity=X_t^\velocity+ \velocity_t(X_t)\in\statespace \label{eq: dynamical system}\,
\end{align}
where $\velocity:(t,x)\in\Nb\x\statespace\to\velocityspace\subset\statespace$. 
Equivalently, letting $\mu_0$ be a probability measure on $(\pspaceomega,\pspacef)$, this can be viewed in terms of the evolution of the law of $X_t$ from its initial distribution $\mu_0$ as
\begin{align}
    \begin{cases} \Lc(X_{t+1}^\velocity)=(\Id+\velocity_t)\pushforward\Lc(X_t^\velocity)\\ \Lc(X_0^\velocity)=\mu_0 \end{cases}
\end{align}

The development of drifting models \cite{deng2026generative} rests on the idea of crafting (by hand) a vector-field-valued map $v$ which is asymptotically stable and whose attracting equilibrium induces $X_{\infty}\sim\nu$ for a given target distribution $\nu\in\Ps(\statespace)$. 
By using the methodology of least-squares fitting outlined in \cite{deng2026generative}, the specification of $v$ can be leveraged into the \textit{training} of a neural network $f$ which yields an approximate pushforward map transforming a given noise distribution $\mu_0$ into the target distribution $\nu$, i.e.\ $\nu\approx f\pushforward\mu_0$.

From a technical standpoint, the authors of \cite{deng2026generative} construct a family of vector fields $v$ depending on $\nu$ and $\Lc(X_{t}^\velocity)$ which admits as \emph{an} equilibrium the target distribution $\nu$. However, under general conditions it remains unclear whether the resulting mean-field dynamics of $X_t$ will meaningfully converge to $\nu$ as $t\to\infty$, despite them giving some sufficient conditions. 

A contrario, it is tempting to reinterpret \eqref{eq: dynamical system} as a control problem, and try and determine $\velocity$ through the specification of an objective function to be optimised. At an intuitive level, we can then view the efforts of \cite{deng2026generative} as an ad-hoc attempt to guess a control $\velocity$ which will yield the desired target distribution $\nu$. 

Inspired by the recent work of \cite{PDRLOT}, we propose to study this as a \emph{measure constrained} optimal control problem.

\section{Discounted setting}

Let $\gamma\in[0,1)$, $c: \statespace^2 \to \mathbb{R}$ be a measurable cost function, and let $\mu_0\in\Ps(\statespace)$ be a given initial distribution, and consider the optimal control problem, 
\begin{align}
    \primalvalue:\mu_0\in\Ps(\statespace)\mapsto\min_{\control\in\controlspace}&\quad \Eb\left[\sum_{t=0}^{\infty} \gamma^t c(X_t^\control, X_{t+1}^\control) \right]\label{eq: optimal control problem}\\
    \mbox{ subject to }& X_{t+1}^\control=X_t^\control+ \velocity_t(X_t^\control) \mbox{ for all } t\ge 0,\notag \\
    &\quad X_0^\control\sim\mu_0,\notag
\end{align}
where $\controlspace$ is a set of admissible controls. By abuse of notation, we will write $\primalvalue(x)$ for the value of the optimal control problem when $\mu_0=\delta_x$ is a Dirac measure at $x\in\statespace$. It is easy to see that a dynamic programming principle holds for $\primalvalue$, namely
\begin{align}
    \primalvalue(x)=\min_{\bar v \in\Sc}\{ c(x, x+\bar v) + \gamma \primalvalue(x+\bar v)\}\,.
\end{align}
and that thus 
\begin{align}
    \primalvalue(\mu_0)=\int \primalvalue(x)\mu_0(\de x)\,.
\end{align}
The optimal policy $\control^*_\cdot\equiv \velocity^*$ is then readily extracted from the minimiser of the right-hand side of the dynamic programming principle.

Consequently, for any $x_0$, optimal trajectories of the system \eqref{eq: dynamical system} follow the constant speed curves
 \begin{align}
    t\mapsto \sum_{s=0}^{t-1} \velocity^*(x_s)= (\Id+v^*)^t x_0,
\end{align} where $x_{s+1}=x_s+\velocity^*(x_s)$.
This suggests an alternative fitting approach than the one proposed in \cite{deng2026generative}: find an approximate fixed point of $\primalvalue$ using least-squares and a Neural Network, somehow extract $v^*$ (for example with $c$ being a squared norm to obtain the network gradient), and then predictions can be computed by feeding forward through the network a few times. However the issue is we don't have a good way to ensure $X_t\to \nu$ from this. But it could give us ideas for after we include the constraint $\Lc(X_t^\velocity)\to\nu$ in the optimal control problem. 

To reduce the simulation cost, we can also try to find an efficient representation of $(\Id + v^*)^t$ for $t$ either fixed or as an input parameter of a neural network approximating it (this gets us closer in spirit to the simulation-free approach of drifting models). If $v^*$ is a gradient field, then we can use the fact that $(\Id + v^*)^t$ is the gradient of a convex potential field perhaps, by writing this as a gradient flow (this would mirror the approach of \cite{PDRLOT}). If $v^*$ is itself a linear operator, then we can directly operate on the spectrum of $v^*$ to compute the exponential for any $t$ very easily. How to justify that $v^*$ is a linear operator is unclear, maybe we can use a latent space / NTK representation or restrict the control class to linear operators. An approximate alternative is presented by linearising $v^*$ around $x_0$ and then computing the exponential of the linearised operator. This corresponds to doing a single shot straight path simulation. I suppose we could also try to be clever by combining different points like in a Runge-Kutta scheme to get very good approximations on the cheap.




\subsection{Formulations of the simulation constraint}

The dynamical system \eqref{eq: dynamical system}, in the equilibrium condition specified by \cite{deng2026generative}, would ideally satify 
\begin{align}
    \lim_{t\to\infty} \Lc(X_t^\velocity)=\nu \mbox{ i.e. } \Lc(X_{t+1}^\velocity)\approx\Lc(X_t^\velocity)\approx\nu\,,
\end{align}
(understood as in distribution) for $t$ sufficiently large, where $\nu\in\Ps(\statespace)$ is a given target distribution, and equality is in distribution. If we restrict ourselves to distributions with finite variance (sledgehammer argument: images are bounded), by the metrisation of weak convergence by the Wasserstein-2 distance, this is equivalent to
\begin{align}
\Wass_2(\Lc(X_t^\velocity),\nu)\approx\Wass_2(\Lc(X_{t+1}^\velocity),\Lc(X_t^\velocity))\approx0\,.
\end{align}
One path: fix as a limit constraint, combine with the other constraint and relax? 
The stationarity condition can be equivalently rewritten in terms of the velocity field $v_t$ as 
\begin{align}
    \sup_{\vp\in C_b(\statespace)}\int\vp(x+v_t(x))- \vp(x)\Lc(X_t^\velocity)(\de x)=0\,,
\end{align}
which can be computed using a GAN but this $\sup$ is very bad. Anticipating a bit, if $\mu^\velocity_t=\Lc(X_t^\velocity,X_{t+1}^\velocity)$ and denoting by $\oplus$ the direct sum, this is a convex constraint on $\mu^\velocity_t$:
\begin{align}
    \mu^\velocity_t\in\Ps_\gamma(\statespace^2): \sup_{\vp\in C_b(\statespace)}\int\vp\oplus\vp(x,y)\mu^\velocity_t(\de x,\de y)=0\,.
\end{align}
Equivalently the distributional constraint can be rewritten as
\begin{align}
    \sup_{\vp\in C_b(\statespace)}\int\vp(x)\nu(\de x)- \vp(x)\Lc(X_t^\velocity)(\de x)=0\,.
\end{align}

One remark about imposing stationarity, if $v$ is constant then $\mu:=\Lc(X_t^\velocity)= \Lc(X_{t+1}^\velocity)$ implies that $\mu=(\Id+v) \pushforward \mu$, and thus $\mu\in \ker(\Id + (\Id+v)\pushforward)$, understood as a linear operator on signed measures. Note that this is a linear constraint on $\mu$ but it is coupled to the primal adjoint variable $v$.

Another way to impose stationarity is to require that $T:=\Id+v$ is admissible as a Monge map for the transport problem between $\mu$ and itself, i.e. $T\pushforward \mu = \mu$. This can perhaps be relaxed to a Kolmogorov coupling constraint on the joint. Come to think of it, this is exactly a rewriting of the constraints that the joint be the joint of the stationary Markov chain with transition $T$. 



\subsection{Dual control}

While the primal theory  tackles \eqref{eq: optimal control problem} head-on, we will instead focus on the (control-theoretic) dual formulation of the problem, which is linear in the space of measures.


To do so, let $\measures_+(\statespace^2)$ be the space of $\sigma$-finite positive measures on $\statespace^2$, and let $\Ps(\statespace^2)\subset \measures(\statespace^2)$ be the space of probability measures on $\statespace^2$.  We define the ``forward'' and ``backward'' marginalisation operators $\fop$ and $\bop$ on $\measures_+(\statespace^2)$ as follows:
\begin{align}
\fop: \mu\in\measures_+(\statespace^2)\to \int_{\statespace} \mu(\de x,\cdot)\in \measures_+(\statespace)\\
\bop: \mu\in\measures_+(\statespace^2)\to \int_{\statespace} \mu(\cdot,\de y)\in \measures_+(\statespace)
\end{align}
For any $\gamma\in[0,1)$, let $\Ps_\gamma(\statespace^2):=(1-\gamma)^{-1}\Ps(\statespace^2)$ be the subspace of $\measures_+(\statespace^2)$ of measures with total mass $(1-\gamma)^{-1}$.

For any $\control\in\controlspace$, with $\mu_0$ implicitely fixed, let $\mu^\control\in\Ps_\gamma(\statespace^2)$ be the measure defined by 
\begin{align}
    \mu^\control(\de x,\de y):=\sum_{t=0}^{\infty} \gamma^t \Lc(X_t^\control,X_{t+1}^\control)(\de x,\de y).
\end{align}
This is equivalently written as \begin{align}
    \mu^\control: (A,B)\in\pspacef^2\to \sum_{t=0}^{\infty} \gamma^t \Pb(X_t^\control\in A, X_{t+1}^\control\in B) = \sum_{t=0}^{\infty} \gamma^t \Pb(X_t^\control\in A,(\Id+\velocity_t)(X_t^\control)\in B)\,.
\end{align}


The \emph{control-theoretic dual} is the linear problem
\begin{align}
\dualvalue:=\min_{\control\in\controlspace} &\quad \int c(x,y)\mu^\control(\de x,\de y) \label{eq: dual control problem}\\
\mbox{ subject to }& \mu^\control\in\Ps_\gamma(\statespace^2),\notag\\
&\quad \bop\mu^\control=\mu_0 + \gamma \fop \mu^\control \\ %\sum_{t=0}^{\infty} \gamma^t \Lc(X_t^\control),
&\quad  \fop\mu^\control= \sum_{t=0}^{\infty} \gamma^t \bigcomp_{s=1}^t(\Id+\velocity_s)\pushforward \mu_0 \notag
\end{align}

\begin{lemma}\label{lemma: control-theoretic duality}
    The control-theoretic duality is strong, \ie, $\primalvalue=\dualvalue$. (We will treat minimiser equivalence later.)
\end{lemma}

\begin{proof}
    That $\int c\de \mu^\control = \Eb[\sum_{t=0}^{\infty} \gamma^t c(X_t^\control, X_{t+1}^\control)]$ is immediate from the definition of $\mu^\control$. 
    
    To establish the constraints of the dual, note that, by definition $\bop\mu^\control=\sum_{t=0}^{\infty} \gamma^t \Lc(X_{t}^\control)$, and thus
    \begin{align}
        \fop\mu^\control=\sum_{t=0}^{\infty} \gamma^t \Lc(X_{t+1}^\control)= \frac1\gamma (\bop\mu^\control - \mu_0). 
    \end{align}
    This establishes the constraint associated to $X_0\sim\mu_0$. The constraint for the dynamics directly leads to 
    \begin{align}
        \fop\mu^\control=\sum_{t=0}^{\infty} \gamma^t \Lc(X_{t+1}^\control)= \sum_{t=0}^{\infty} \gamma^t (\Id+\velocity_t)\pushforward\Lc(X_t^\control) = \sum_{t=0}^{\infty} \gamma^t \bigcomp_{s=1}^t(\Id+\velocity_s)\pushforward \mu_0\,.
    \end{align}
\end{proof}


\subsection{A Remark on minimal energy configurations}

If we set $c(x,y)=\|x-y\|^2$, then the optimal control problem \eqref{eq: optimal control problem} becomes 
\begin{align}
    \min_{\control\in\controlspace}&\quad \sum_{t=0}^\infty \gamma^t \norm{\velocity_t}^2_2\\
    \mbox{ subject to }& X_{t+1}^\control=X_t^\control+ \velocity_t(X_t) \mbox{ for all } t\ge 0,\notag \\
    &\quad X_0\sim\mu_0,\notag
\end{align}

The constraints vanish in this case, as the objective is independent of the distribution of $X_t^\velocity$. The optimal control problem then reduces to a standard optimal control problem, and the optimal control is given by $\velocity^*_t\equiv 0$ for all $t\ge 0$.

For this to be useful for the transport problem, we need to balance with the constraints so that the movement is uniform in time rather than concentrated at infinity.


\section{Finite Horizon formulations}

If we set the problem as finite horizon, it is much easier to lay out the constraint we care about for generation.

In particular, we can consider the following finite-horizon optimal control problem, in continuous time, with a terminal constraint:
\begin{align}
    \primalvalue(0,\mu_0):=\min_{\control\in\controlspace_0}&\quad \int\int_0^1 c(X_t^\control,\velocity_t(X_t^\control))\de t\de \mu_0\notag\\
    \mbox{ subject to }& \dot X_t^\control=\velocity_t(X_t^\control) \mbox{ for all } t\in[0,1],\\
    &\quad X_0^\control\sim\mu_0,\notag\\
    &\quad \Lc(X_1^\control)=\nu\notag
\end{align}
We can then define the cost-to-go functions as $\primalvalue(t,\mu):=\min_{\control\in\controlspace_t}\int_t^1 c(X_s^\control,\velocity_s(X_s^\control))\de s$ subject to $\dot X_s^\control=\velocity_s(X_s^\control)$ for all $s\in[t,1]$, $X_t^\control\sim\mu$, and $\Lc(X_1^\control)=\nu$.
In either case $\controlspace_t$ is the set of admissible controls on $[t,1]$. 

When $c(x,v)=\norm{v}^2_2$, this is exactly the Benamou-Brenier formulation of optimal transport
\begin{align}
\Wass_2^2(\mu_0,\nu)=\min_{\control\in\controlspace_0}&\quad \int_0^1 \int \norm{\velocity_t(x)}^2_2\de \mu_t(x)\de t\\
    \mbox{ subject to }& \frac\de{\de t}\mu_t+\nabla\cdot(\velocity_t\mu_t)=0 \mbox{ for all } t\in[0,1],\\
    &\quad \mu_0=\mu_0,\notag\\
    &\quad \mu_1=\nu\notag\\
    &\quad t\mapsto\mu_t \mbox{ a.c. and } \Ps_2(\statespace).
\end{align} 


Discretising the problem as in \citet{PDRLOT}, we can solve this via a primal-dual approach, including with only sample access to $\nu$. 

This forgets about the stationarity completely, but notice that any convex function of $\velocity$ induces a stationary control and thus immediate stationarity of the velocity. 

Something that could be interesting is to include in the cost a term such as $\norm{x}_1$ to encourage a turnpike behaviour, pushing into a latent space and back out of it. We can also replace the squared norm with $s\norm{v}_2^2$ or with $\norm{v}_2$ in order to encourage faster convergence to the target distribution from $\mu_0$. 


\section{Alternative mean-field control formulations}

Another way to think about the drifting model optimisation is to treat the empirical data as a finite-particle approximation of a mean-field control problem. In particular, we can view \eqref{eq: optimal control problem} as the mean-field control problem
\begin{align}
    \primalvalue(\mu_0):=\min_{\control\in\controlspace}&\quad \Eb\left[\sum_{t=0}^{\infty} \gamma^t c(X_t^\control, \Lc(X_t^\control),\velocity_t) \right]\\
    \mbox{ subject to }& X_{t+1}^\control=X_t^\control+ \velocity_t(X_t^\control) \mbox{ for all } t\ge 0,\notag \\
    &\quad X_0^\control\sim\mu_0,\notag
\end{align}
where $c$ is now a function of the law of $X_t^\control$ and the velocity field $\velocity_t$. If we take, for instance, $c(x,\mu,v)=\norm{v}^2_2+\Wass_2^2(\mu,\nu)$, then the optimal control problem becomes
\begin{align}
    \primalvalue(\mu_0):=\min_{\control\in\controlspace}&\quad \Eb\left[\sum_{t=0}^{\infty} \gamma^t \norm{\velocity_t(X_t^\control)}^2_2+\Wass_2^2(\Lc(X_t^\control),\nu) \right]\\
    \mbox{ subject to }& X_{t+1}^\control=X_t^\control+ \velocity_t(X_t^\control,\Lc(X_t^\control)) \mbox{ for all } t\ge 0,\notag \\
    &\quad X_0^\control\sim\mu_0,\notag
\end{align}
which is (we can always change the decay scheme of costs to have finite horizon) a penalised version of the Benamou-Brenier tranport problem which encourages convergence to the target distribution $\nu$. Of course, the downside of this is that the velocity field will depend on the attraction of all particles to all samples, just as in the original drifting model. 


Below is a small generated overview of what we can hope to do in terms of resolution. Basically, everything will work the same way as with non-measure-valued states, but there will be a nastier measure-valued Bellman equation \eqref{eq: mf master} to solve. However, the resulting optimal control is still a feedback control, and there is a rich literature on the resolution using, say JKO-type schemes. The path to samples it to view our algorithm as directly being the numerical scheme on this problem, including a sample collected to represent the two measures $\mu$ and $\nu$ via propagation of chaos. 


Restricting to Markov feedback controls $\velocity_t(\cdot)=\alpha(\cdot,\mu_t)$ with $\mu_t=\Lc(X_t)$, the induced one-step map is
\[
    T^\alpha_\mu(x):=x+\alpha(x,\mu),\qquad \mu^+=(T^\alpha_\mu)\pushforward\mu.
\]
The dynamic programming principle on $\Ps_2(\statespace)$ then takes the form
\begin{align}
    \primalvalue(\mu)=\inf_{\alpha}\left\{\int_{\statespace} c(x,\mu,\alpha(x,\mu))\,\mu(\de x)+\gamma\primalvalue\big((T^\alpha_\mu)\pushforward\mu\big)\right\}.\label{eq: mf dpp}
\end{align}
This is the natural discrete-time Bellman equation on the Wasserstein space of probability measures; see for instance \citet{PW18,BFY13,CD18I}. In particular, the optimal feedback should be characterised by minimisers of the right-hand side of \eqref{eq: mf dpp}, exactly as in the finite-dimensional case, but now after lifting the state variable from $x$ to the whole law $\mu$.

Assuming that $\primalvalue$ is differentiable in the sense of Lions\footnote{See e.g. \citet{bao2021bismut}, basically it's Fréchet derivative but for measures. We can easily assume densities and work with them instead.}, one may define the decoupling field
\[
    U(\mu,x):=\frac{\delta \primalvalue}{\delta\mu}(\mu)(x),
\]
and then formally differentiate \eqref{eq: mf dpp} with respect to $\mu$. If $\alpha^*(\cdot,\mu)$ is an optimal feedback and $T^*_{\mu}(x):=x+\alpha^*(x,\mu)$ with $\mu^+=(T^*_{\mu})\pushforward\mu$, this leads to the discrete-time master equation
\begin{align}
    U(\mu,x)=&\ c(x,\mu,\alpha^*(x,\mu))+\gamma U(\mu^+,T^*_{\mu}(x))+\int_{\statespace}\frac{\delta c}{\delta\mu}(y,\mu,\alpha^*(y,\mu))(x)\,\mu(\de y)\label{eq: mf master}\\
    &\ +\gamma\int_{\statespace} D_xU(\mu^+,T^*_{\mu}(y))\cdot \frac{\delta T^*_{\mu}}{\delta\mu}(y)(x)\,\mu(\de y),\notag
\end{align}
again at a purely formal level. Equation \eqref{eq: mf master} is the discrete analogue of the continuous-time master equation from mean-field game and mean-field type control theory, and should be understood as a fixed-point equation for the Lions derivative of the value functional; compare with \citet{CDLL19,CD18II}. It may be possible to formulate this more cleanly by first passing to a finite-horizon problem and then sending the horizon to infinity.

As for solving the problem, the quadratic term $\norm{v}^2_2$ by itself does not make the optimisation a simple least-squares problem: the nonlinearity is hidden in the pushforward constraint $\mu^+=(\mathrm{id}+\alpha(\cdot,\mu))\pushforward\mu$ and in the dependence of the running cost on the evolving law. Thus a direct calculus-of-variations approach is plausible only in special regimes where one can relax the control problem to a convex optimisation over curves of measures or transport plans, as in Benamou--Brenier or JKO-type formulations; see \citet{San15,JKO98,AGS05}. In the general mean-field control setting one typically needs either the dynamic programming/master equation route or the stochastic/variational maximum principle developed in \citet{BFY13,CD18I,CD18II}. 

\section{Transient distribution model}

The original drifting model paper \cite{deng2026generative} considers a Markov chain over the training time with the goal of following its transition kernel to eventually reach a stationary distribution which is close to or equal to the target distribution $\nu$. In the paper, the transition kernel is constructed by hand, with some conditions which should ensure that the stationary distribution is indeed $\nu$. 

In contrast, it is tempting to view the problem as a control problem, and to explicitely optimise the transition kernel to reach the target distribution $\nu$ in as few learning steps as possible. This imposes a constraint on the ergodic control problem induced by the controlled Markov chain induced by all possible vector fields $v$ during the learning dynamics. To handle this difficulty, we propose to analyse the resulting problem directly through a Lagrangian relaxation with the goal of performing a primal-dual optimisation over the vector field $v$ in which we will use samples for access to $\nu$. 

In practice, much like in the original drifting model paper, we will subsequently use optimise over a parametric family of vector fields as we will parametrise the generative model by a neural network. Because of the chain rule, we will need to compute the gradients of the full Lagrangian with respect to $v$, which is the paradigm we adopt here to begin with. Note that, using the trick of \cite{deng2026generative}, we don't need to explicitely differentiate through the network, it is sufficient to find a good $v$ and then minimise the least-squares loss to stationarity, or use the formulation of \cite{PDRLOT} and generate directly from samples of $\mu_0$ via the transport $v$ (setting the initial network to the identity).

\subsection{Lagrangian of the problem}

If we let $v$ be a constant velocity field, assuming that $v$ induces a stationary distribution $q_v$ for the Markov chain \eqref{eq: dynamical system}, i.e. $q_v=(\Id+v)\pushforward q_v$, then we can define the transient distribution of the Markov chain as
\begin{align}
    \phi:A\in\Bs(\statespace)\to \sum_{t=0}^\infty \Pb(X_t^\velocity\in A)-q_v(A)\in[0,\infty].
\end{align}

In this case, we can aim to optimise over $v$ to minimise the energy of the transient distribution, according to some cost criterion. For instance, we can consider the following optimisation problem:
\begin{align}
    \min_{v\in\Sc}& \int_{\statespace} c(x)\phi(\de x)\\ \mbox{ subject to }& q_v=(\Id+v)\pushforward q_v.
\end{align}

\begin{remark}
    It's important to note how different this formulation is from the Value Driven Transport paper of \cite{PDRLOT}. There is no transport cost from $x$ to $x+v(x)$, here (at least not explicitely). What is penalised is a weighting of the system's convergence rate to the stationary distribution based on different costs in space. If we proceed we should be able to show this effect by making $c=1+\1_{S}$ for a set $S$ of states, such as a halfspace. We should in general take $c\equiv 1$ if we have no prior knowledge of the target distribution $\nu$. 
\end{remark}


A simple calculation shows that 
\begin{align}
    \phi_v=\mu_0-q_v+ (\Id+ v)\pushforward \phi_v, \label{eq: transient distribution constraint}
\end{align}
if and only if $\phi_v$ is defined as above. We can add this linear constraint to the optimisation problem, and relax the objective to separate $\phi_v$ and $v$. This gives
\begin{align}
    \min_{v\in\Sc,\phi\in\Ms}& \langle c\vert\phi\rangle\\ \mbox{ subject to }& q_v=(\Id+v)\pushforward q_v,\\ & \phi_v=\mu_0-q_v+ (\Id+ v)\pushforward \phi_v \\& \phi\ge \phi_v.
\end{align}

The transient distribution $\phi_v$ also has two properties which we can switch to constraints on $\phi$ to make optimising over $\phi$ easier. First, $\phi(\Omega)=0$, as both $\Pb(X_t^\velocity\in\cdot)$ and $q_v$ are probability measures. Notice that $\phi$ is unbounded, as is easily shown by taking the two-state system with transition probability matrix
\begin{align}
    P_v=\begin{pmatrix} 1-\ve & \ve\\ 0 & 1\end{pmatrix}
\end{align}
wherein the transient part has mass $(1-\ve)^{-1}$, on $s_1$ and $1-(1-\ve)^{-1}$ on $s_2$, while $q(s_2)=1$ for any $\ve>0$. Sending $\ve\to 1$ sends $\phi(s_1)\to\infty$ and $\phi(s_2)\to -\infty$.

Adding this, plus the constraint that $q_v$ be the target distribution $\nu$, we can write the optimisation problem as
\begin{align}
    \min_{v\in\Sc,\phi\in\Ms}& \langle c\vert\phi\rangle\\ \mbox{ subject to }& q_v=(\Id+v)\pushforward q_v,\\ & \phi_v=\mu_0-q_v+ (\Id+ v)\pushforward \phi_v\\ & q_v=\nu \\& \phi\ge \phi_v,\\ & \phi(\Omega)=0.
\end{align}

While this system has non-linearities (in $v$) in the constraints, it doesn't prevent us from writing the primal-dual system and trying to solve it. The main issue for solving is that we need to be able to 1) compute the stationary distribution $q_v$ for any $v$,  2) compute the transient distribution $\phi_v$ for any $v$, and 3) compute the gradient of the transient distribution with respect to $v$ (maybe doable via PG theorem ?). 

The Lagrangian for the system is 
\begin{align}
    \Lc(v,\phi,\lambda_1,\lambda_2,\lambda_3,\lambda_4)=&\langle c\vert\phi\rangle+\langle \lambda_1\vert q_v-(\Id+v)\pushforward q_v\rangle+\langle \lambda_2\vert \phi_v-\mu_0+q_v-(\Id+v)\pushforward \phi_v\rangle\notag\\
    &+\langle \lambda_3\vert q_v-\nu\rangle+\langle \lambda_4\vert \phi_v-\phi\rangle+\lambda_5\phi(\Omega)\,,
\end{align}
or, reordered,
\begin{align}
    \Lc(v,\phi,\lambda_1,\lambda_2,\lambda_3,\lambda_4)=&\langle c-\lambda_4\vert\phi\rangle+\lambda_5\phi(\Omega) + \langle \lambda_3\vert\nu\rangle+ \langle  \lambda_1+\lambda_2+\lambda_3\vert q_v\rangle+ \langle \lambda_2+\lambda_4\vert \phi_v\rangle \notag\\
    &\quad- \langle \lambda_1 \vert P_vq_v\rangle -\langle \lambda_2 \vert P_v\phi_v\rangle\,,\label{eq: lagrangian with measures}
\end{align}
which is linear in $\phi$ and $q_v$, but non-linear in $v$.

\subsection{Gradients of the Lagrangian}

The gradients with respect to the multipliers and $\phi$ are easy provided knowledge of $\phi_v$ and $q_v$ (except for $\lambda_3$ which requires the target $\nu$ but it can take SGD). On the other hand, taking the gradient in the control space yields\footnote{For simplicity, we develop below for $c$ independent of $v$, but the same reasoning applies if $c$ depends on $v$, up to adding $\langle \nabla_v c\vert \phi\rangle$.}
\begin{align}
    \nabla_v \Lc(v,\phi,\lambda_1,\lambda_2,\lambda_3,\lambda_4)=&\nabla_v\left(\langle  \lambda_1+\lambda_2+\lambda_3\vert q_v\rangle+ \langle \lambda_2+\lambda_4\vert \phi_v\rangle\right)\notag\\
    &\qquad-\nabla_v\left(\langle \lambda_1 \vert P_vq_v\rangle +\langle \lambda_2 \vert P_v\phi_v\rangle\right) 
\end{align}
where $P_v:=(\Id+v)\pushforward$ is the transition operator of the Markov chain, viewed as a linear operator on signed measures.

We need to be a bit careful about topology here in order to get these derivatives into the inner products. So to be general, let's define a function class $\Fs$ of functions, which is included in $\Cc^0_b$ (the pre-dual of $\Ms$ the space of finite signed measures) and the corresponding $\Fs$-weak topology on the space of signed measures which is the smallest topology such that $\mu\mapsto \langle f\vert \mu\rangle$ is continuous for all $f\in\Fs$. (In general, we will want to take a dense subset of $\Cc^0_b$ with nice test-function properties, such as $\Cc^\infty_c$ or $\Cc^\infty_0$.)

\begin{definition}[Gateaux derivative in the $\Fs$-weak topology]
    Let $\Fs\subset \Cc^0_b(\statespace)$ be a function class, and let $v\mapsto m_v$ be a map from $\Sc$ to $\Ms$. We say that $m_v$ is $\Fs$-weakly Gateaux differentiable at $v\in\Sc$ if there exists $\De_h m_v:\Sc\to \Ms\in\Ms$ such that, 
    \begin{align}
        \lim_{\varepsilon\to 0} \sup_{f\in\Fs}\left|\frac{\langle f\vert m_{v+\varepsilon h}-m_v\rangle}{\varepsilon}-\langle f\vert \De_h m_v\rangle\right|=0.
    \end{align}
    We call $\De_h m_v$ the $\Fs$-weak Gateaux derivative of $m_v$ at $v$ in direction $h\in\Sc$, and, if $h\mapsto \De_h m_v$ is well-defined for all $v$ we call $v\mapsto \De_h m_v$ the Gateaux derivative of $m_v$ in direction $h$.
\end{definition}

Similarly, we can define the Gateaux derivative of an operator $v\mapsto O_v$ with values in linear operators on signed measures, by calling $\bar\De_h O_v$ the Gateaux derivative of $O_v$ at $v$ in direction $h$ if, for all $\mu\in\Ms$,
\begin{align}
    \lim_{\varepsilon\to 0} \sup_{f\in\Fs}\left|\frac{\langle f\vert O_{v+\varepsilon h}\mu-O_v\mu\rangle}{\varepsilon}-\langle f\vert \bar\De_h O_v\mu\rangle\right|=0.
\end{align}
(We use the notation $\bar\De$ to distinguish the derivative of an operator applied to a measure-valued map from the derivative of the composed map $v\mapsto O_v m_v$ itself.)
Assuming both $v\mapsto m_v$ and $v\mapsto O_v$ are Gateaux differentiable, we can then write the product rule
\begin{align}
    \De_h (O_v m_v)=\bar\De_h O_v m_v+O_v \De_h m_v. \label{eq: product rule for operator and measure}
\end{align}
Thus, we can rewrite \eqref{eq: lagrangian with measures} as
\begin{align}
    \nabla_v \Lc(v,\phi,\lambda_1,\lambda_2,\lambda_3,\lambda_4)=&\langle \lambda_1+\lambda_2+\lambda_3\vert \De_h q_v\rangle+ \langle \lambda_2+\lambda_4\vert \De_h \phi_v\rangle\notag\\
    &\qquad-\langle \lambda_1 \vert \bar\De_h P_v q_v+ P_v \De_h q_v\rangle -\langle \lambda_2 \vert \bar\De_h P_v \phi_v+ P_v \De_h \phi_v\rangle. \label{eq: lagrangian gradient with measures}
\end{align}

\begin{lemma}\label{lemma: derivative of pushforward operator}
    Let $\mu\in\Ms$ and $f\in\Fs$. Assume $\Fs\subset\Cc^1(\statespace)$ and then, for any $h\in\Sc$ and $v\in\Sc$, the duality pairing of $f$ and the  Gateaux derivative of the pushforward operator $P_v$ is given by
    \begin{align}
        \langle f\vert \bar\De_h P_v \mu\rangle=\int \nabla f(x+v(x))\cdot h(x)\,\mu(\de x).
    \end{align}
    Which, conceptually can be understood as operator on signed measures whose adjoint is $(\bar\De_h P_v)^*:f \mapsto h\nabla f(\cdot+v(\cdot)) $, where the gradient is of course only on the argument of $f$ (is this what kids these days call a stopgradient?).

    Consequently, $v\mapsto \langle f\vert P_v \mu\rangle$ is Fréchet differentiable in the $\Fs$-weak topology, with Fréchet derivative
    \begin{align}
        h\mapsto \int h(x)\nabla f(x+v(x))\,\mu(\de x).
    \end{align}
\end{lemma}

\begin{proof}
    This follows by a taylor expansion of $f$ and the definition of the pushforward operator:
    \begin{align}
        \langle f\vert P_{v+\varepsilon h}\mu\rangle&=\int f(x+v(x)+\varepsilon h(x))\,\mu(\de x)\\&=\int f(x+v(x))\,\mu(\de x)+\varepsilon \int \nabla f(x+v(x))\cdot h(x)\,\mu(\de x)+o(\varepsilon),
    \end{align}
    and thus
    \begin{align}
        \langle f\vert \bar\De_h P_v \mu\rangle=\lim_{\varepsilon\to 0} \frac{\langle f\vert P_{v+\varepsilon h}\mu-P_v\mu\rangle}{\varepsilon}=\int \nabla f(x+v(x))\cdot h(x)\,\mu(\de x).
    \end{align}
\end{proof}

While for the $P_v\mu$ terms in the lagrangian ($\mu\in\{q_v,\phi_v\}$), it was convenient to differentiate the inner product with a test function $f\in\Fs$ and directly express the differential of the pushforward operator in terms of the adjoint action on $f$, when differentiating the ergodic terms ($q_v$ and $\phi_v$), we will need to face the measures themselves, and thus we need to compute the Gateaux derivatives $\De_h q_v$ and $\De_h \phi_v$, which we do in \cref{lemma: derivative of q_v} and \cref{lemma: derivative of phi_v} below.

Before we move onto these rigourous constructions, let us recall that we care about differentiative $v\mapsto \langle f\vert \mu\rangle$, in other words, we want to prove a Policy Gradient Theorem of the following form.

\begin{theorem}\label{thm: policy gradient theorem}
    Let $q_v$ be the stationary distribution associated with the operator $P_v$. Let $f\in\Fs$ be a test function, then $v\mapsto \langle f\vert q_v\rangle$ is $\Fs$-weakly Gateaux differentiable at $v$, and its Gateaux derivative in direction $h\in\Sc$ is given by
    \begin{align}
        \De_h \langle f\vert q_v\rangle=\int \nabla u_f(x+v(x))\cdot h(x)\,q_v(\de x),
    \end{align}
    where $u_f$ is a solution to the Poisson equation
    \begin{align}
        u_f-P_v^*u_f=f-\langle f\vert q_v\rangle.
    \end{align}
\end{theorem}

\begin{proof}
    Under their technical assumptions, the result follows directly from \cref{lemma: derivative of q_v} and \cref{lemma: poisson representation of derivative of q_v}.  
\end{proof}

\begin{remark}
    Most often, policy gradient theorems are expressed in terms of the derivative of the score (log-transform of the transition density) of a softmax policy, this will generally not be possible in our setting, as $P_v$ is a transport operator.  
\end{remark}

Let us now compute the Gateaux derivative of the stationary distribution $q_v$ with respect to $v$ to prove \cref{thm: policy gradient theorem}. 


\begin{lemma}\label{lemma: derivative of q_v}
    Let $q_v$ be the stationary distribution associated with the operator $P_v$. Assume that for all sufficiently small $\varepsilon>0$, $(v+\varepsilon h\in\mathcal S)$, the operator $(P_{v+\varepsilon h})$ is a Markov operator on probabilities, and it admits a stationary distribution $(q_{v+\varepsilon h})$.\footnote{We need to be a bit careful about verifying these assumptions with stability properties.}
    
    Assume that $v\mapsto q_v$ is $\Fs$-weakly Gateaux differentiable at $v$, and that $(\Id-P_v)$ admits an inverse on the subspace of zero-mass signed measures. 
    
    Then, for any $h\in\Sc$, the Gateaux derivative of $q_v$ in direction $h$ is the unique zero-mass signed measure $m$ solving
    \begin{align}
        (\Id-P_v) m =\bar\De_h P_v q_v,
    \end{align}
    (understood in the $\Fs$-weak topology, still), that is,
    \begin{align}
        \De_h q_v=(\Id-P_v)^{-1}\bar\De_h P_v q_v.
    \end{align}
    If, in addition, the Neumann series converges on zero-mass signed measures, then
    \begin{align}
        \De_h q_v=\sum_{t=0}^\infty P_v^t\bar\De_h P_v q_v.
    \end{align}
\end{lemma}

\begin{proof}
    Since $q_v$ is stationary, it satisfies
    \begin{align}
        q_v=P_v q_v.
    \end{align}
    Differentiating this identity in the direction $h$ and applying the product rule established in \eqref{eq: product rule for operator and measure} gives
    \begin{align}
        \De_h q_v=\De_h(P_v q_v)=\bar\De_h P_v q_v+P_v\De_h q_v.
    \end{align}
    Rearranging yields
    \begin{align}
        (\Id-P_v)\De_h q_v=\bar\De_h P_v q_v.
    \end{align}
    Moreover, since $q_{v+\varepsilon h}$ is a probability measure for every sufficiently small $\varepsilon$, we also have
    \begin{align}
        \De_h q_v(\Omega)=\lim_{\varepsilon\to 0}\frac{q_{v+\varepsilon h}(\Omega)-q_v(\Omega)}{\varepsilon}=0.
    \end{align}
    Thus $\De_h q_v$ belongs to the zero-mass subspace, on which $(\Id-P_v)$ is invertible by assumption, so
    \begin{align}
        \De_h q_v=(\Id-P_v)^{-1}\bar\De_h P_v q_v.
    \end{align}
    If the Neumann series for $(\Id-P_v)^{-1}$ converges on that subspace, then
    \begin{align}
        (\Id-P_v)^{-1}=\sum_{t=0}^\infty P_v^t,
    \end{align}
    and therefore
    \begin{align}
        \De_h q_v=\sum_{t=0}^\infty P_v^t\bar\De_h P_v q_v.
    \end{align}
\end{proof}

Finally, let's rewrite the duality bracket $\langle f\vert \De_h q_v\rangle$ in terms of the solution to the Poisson equation, which will give us the form of the policy gradient theorem.

\begin{lemma}\label{lemma: poisson representation of derivative of q_v}
    Let $f\in\Fs$, and assume that the Poisson equation
    \begin{align}
        u_f-P_v^*u_f=f-\langle f\vert q_v\rangle
    \end{align}
    admits a solution $u_f\in\Cc^1(\statespace)\cap\Fs$. Then, for every direction $h\in\Sc$,
    \begin{align}
        \De_h\langle f\vert q_v\rangle=\langle u_f\vert \bar\De_h P_v q_v\rangle.
    \end{align}
    Consequently, 
    \begin{align}
        \De_h\langle f\vert q_v\rangle=\int \nabla u_f(x+v(x))\cdot h(x)\,q_v(\de x).
    \end{align}
\end{lemma}

\begin{proof}
    Since $\De_h q_v$ has zero total mass by \cref{lemma: derivative of q_v}, we have
    \begin{align}
        \De_h\langle f\vert q_v\rangle
        =\langle f\vert \De_h q_v\rangle
        =\left\langle f-\langle f\vert q_v\rangle\middle\vert \De_h q_v\right\rangle.
    \end{align}
    Using the Poisson equation, this becomes
    \begin{align}
        \De_h\langle f\vert q_v\rangle
        =\langle (\Id-P_v^*)u_f\vert \De_h q_v\rangle
        =\langle u_f\vert (\Id-P_v)\De_h q_v\rangle.
    \end{align}
    Applying \cref{lemma: derivative of q_v} yields
    \begin{align}
        \De_h\langle f\vert q_v\rangle=\langle u_f\vert \bar\De_h P_v q_v\rangle.
    \end{align}
    The final identity then follows from \cref{lemma: derivative of pushforward operator}:
    \begin{align}
        \langle u_f\vert \bar\De_h P_v q_v\rangle
        =\int \nabla u_f(x+v(x))\cdot h(x)\,q_v(\de x).
    \end{align}
\end{proof}


\begin{remark}
    The existence of a $\Cc^1$ solution to the Poisson equation for any sufficiently regular source term $f$ is a regularity theory result for the associated integral equation. It should be well-studied but will require assumptions. 
\end{remark}

We can proceed in much the same way for the transient distribution $\phi_v$. However, its implicit dependence on $q_v$ (see \eqref{eq: transient distribution constraint}) will reintroduce the derivative of $q_v$ to the calculations. This produces a second coupled Poisson equation on the advantage function.

\begin{lemma}\label{lemma: derivative of phi_v}
    Assume that $v\mapsto \phi_v$ is $\Fs$-weakly Gateaux differentiable at $v$, and that $(\Id-P_v)$ admits an inverse on the subspace of zero-mass signed measures. Then, for any $h\in\Sc$, the Gateaux derivative of $\phi_v$ in direction $h$ is the unique zero-mass signed measure solving
    \begin{align}
        (\Id-P_v)\De_h\phi_v=-\De_h q_v+\bar\De_h P_v\phi_v,
    \end{align}
    that is,
    \begin{align}
        \De_h\phi_v=(\Id-P_v)^{-1}\left(-\De_h q_v+\bar\De_h P_v\phi_v\right).
    \end{align}
\end{lemma}

\begin{proof}
    Differentiating the fixed-point identity
    \begin{align}
        \phi_v=\mu_0-q_v+P_v\phi_v
    \end{align}
    in the direction $h$ and using \eqref{eq: product rule for operator and measure} gives
    \begin{align}
        \De_h\phi_v=-\De_h q_v+\bar\De_h P_v\phi_v+P_v\De_h\phi_v.
    \end{align}
    Hence
    \begin{align}
        (\Id-P_v)\De_h\phi_v=-\De_h q_v+\bar\De_h P_v\phi_v.
    \end{align}
    Since $\phi_v(\Omega)=0$ for every admissible $v$, we also have $\De_h\phi_v(\Omega)=0$, so $\De_h\phi_v$ belongs to the zero-mass subspace. The claimed representation follows by inverting $(\Id-P_v)$ on that subspace.
\end{proof}

\begin{theorem}[Transient measure gradient theorem]\label{thm: transient policy gradient theorem}
    Let $f\in\Fs$, and assume that the Poisson equations
    \begin{align}
        u_f-P_v^*u_f&=f-\langle f\vert q_v\rangle,\\
        w_f-P_v^*w_f&=u_f-\langle u_f\vert q_v\rangle
    \end{align}
    admit solutions $u_f,w_f\in\Cc^1(\statespace)\cap\Fs$. Then the map $v\mapsto \langle f\vert \phi_v\rangle$ is $\Fs$-weakly Gateaux differentiable at $v$, and for every direction $h\in\Sc$,
    \begin{align}
        \De_h\langle f\vert \phi_v\rangle
        =\langle u_f\vert \bar\De_h P_v\phi_v\rangle-\langle w_f\vert \bar\De_h P_v q_v\rangle.
    \end{align}
    Consequently,
    \begin{align}
        \De_h\langle f\vert \phi_v\rangle
        =\int \nabla u_f(x+v(x))\cdot h(x)\,\phi_v(\de x)
        -\int \nabla w_f(x+v(x))\cdot h(x)\,q_v(\de x).
    \end{align}
\end{theorem}

\begin{proof}
    Since $\phi_v(\Omega)=0$ for every admissible $v$, the derivative also has zero total mass, and hence
    \begin{align}
        \De_h\langle f\vert \phi_v\rangle
        =\langle f\vert \De_h\phi_v\rangle
        =\left\langle f-\langle f\vert q_v\rangle\middle\vert \De_h\phi_v\right\rangle.
    \end{align}
    Using the Poisson equation for $u_f$ gives
    \begin{align}
        \De_h\langle f\vert \phi_v\rangle
        =\langle (\Id-P_v^*)u_f\vert \De_h\phi_v\rangle
        =\langle u_f\vert (\Id-P_v)\De_h\phi_v\rangle.
    \end{align}
    Applying \cref{lemma: derivative of phi_v} yields
    \begin{align}
        \De_h\langle f\vert \phi_v\rangle
        =-\langle u_f\vert \De_h q_v\rangle+\langle u_f\vert \bar\De_h P_v\phi_v\rangle. \label{eq: derivative of IP with transient distribution proof}
    \end{align}
    Since $\De_h q_v(\Omega)=0$, we may subtract the constant $\langle u_f\vert q_v\rangle$ and use the Poisson equation for $w_f$ to write
    \begin{align}
        \langle u_f\vert \De_h q_v\rangle
        &=\left\langle u_f-\langle u_f\vert q_v\rangle\middle\vert \De_h q_v\right\rangle\\
        &=\langle (\Id-P_v^*)w_f\vert \De_h q_v\rangle\\
        &=\langle w_f\vert (\Id-P_v)\De_h q_v\rangle.
    \end{align}
    Using \cref{lemma: derivative of q_v} then gives
    \begin{align}
        \langle u_f\vert \De_h q_v\rangle=\langle w_f\vert \bar\De_h P_v q_v\rangle.
    \end{align}
    Substituting back into \eqref{eq: derivative of IP with transient distribution proof}, we obtain
    \begin{align}
        \De_h\langle f\vert \phi_v\rangle
        =\langle u_f\vert \bar\De_h P_v\phi_v\rangle-\langle w_f\vert \bar\De_h P_v q_v\rangle.
    \end{align}
    The integral representation follows from \cref{lemma: derivative of pushforward operator} applied to both terms.
\end{proof}

We can now substitute the stationary and transient policy gradient formulas into the directional derivative of the Lagrangian.

\begin{lemma}\label{lemma: lagrangian gradient via poisson equations}
    Let
    \begin{align}
        f_q&:=(\Id-P_v^*)\lambda_1+\lambda_2+\lambda_3,\\
        f_\phi&:=(\Id-P_v^*)\lambda_2+\lambda_4.
    \end{align}
    Assume that the Poisson equations
    \begin{align}
        U_q-P_v^*U_q&=f_q-\langle f_q\vert q_v\rangle,\\
        U_\phi-P_v^*U_\phi&=f_\phi-\langle f_\phi\vert q_v\rangle,\\
        W_\phi-P_v^*W_\phi&=U_\phi-\langle U_\phi\vert q_v\rangle
    \end{align}
    admit solutions in $\Cc^1(\statespace)\cap\Fs$. Then the directional derivative of the Lagrangian in the direction $h\in\Sc$ is given by
    \begin{align}
        \De_h\Lc(v,\phi,\lambda_1,\lambda_2,\lambda_3,\lambda_4)
        =&\ \langle U_\phi-\lambda_2\vert \bar\De_h P_v\phi_v\rangle\notag\\
        &+\langle U_q-W_\phi-\lambda_1\vert \bar\De_h P_v q_v\rangle.
    \end{align}
    Equivalently,
    \begin{align}
        \De_h\Lc(v,\phi,\lambda_1,\lambda_2,\lambda_3,\lambda_4)
        =&\int \nabla (U_\phi-\lambda_2)(x+v(x))\cdot h(x)\,\phi_v(\de x)\notag\\
        &+\int \nabla (U_q-W_\phi-\lambda_1)(x+v(x))\cdot h(x)\,q_v(\de x).
    \end{align}
    In other words, when this expression is well-defined, the Fréchet derivative of the Lagrangian in the $\Fs$-weak topology is given by
    \begin{align}
        \nabla_v \Lc(v,\phi,\lambda_1,\lambda_2,\lambda_3,\lambda_4)=P_v^*\nabla (U_\phi-\lambda_2)\de\phi_v+P_v^*\nabla (U_q-W_\phi-\lambda_1)\de q_v.
    \end{align}
\end{lemma}

\begin{proof}
    Starting from the directional derivative formula in \eqref{eq: lagrangian gradient with measures} and moving $P_v$ to the test-function side, we obtain
    \begin{align}
        \De_h\Lc(v,\phi,\lambda_1,\lambda_2,\lambda_3,\lambda_4)
        =&\left\langle (\Id-P_v^*)\lambda_1+\lambda_2+\lambda_3\middle\vert \De_h q_v\right\rangle\notag\\
        &+\left\langle (\Id-P_v^*)\lambda_2+\lambda_4\middle\vert \De_h \phi_v\right\rangle\notag\\
        &-\langle \lambda_1\vert \bar\De_h P_v q_v\rangle-\langle \lambda_2\vert \bar\De_h P_v\phi_v\rangle.
    \end{align}
    By definition of $f_q$ and $f_\phi$, this becomes
    \begin{align}
        \De_h\Lc(v,\phi,\lambda_1,\lambda_2,\lambda_3,\lambda_4)
        =&\ \De_h\langle f_q\vert q_v\rangle+\De_h\langle f_\phi\vert \phi_v\rangle\notag\\
        &-\langle \lambda_1\vert \bar\De_h P_v q_v\rangle-\langle \lambda_2\vert \bar\De_h P_v\phi_v\rangle.
    \end{align}
    Applying \cref{thm: policy gradient theorem,thm: transient policy gradient theorem} gives
    \begin{align}
        \De_h\Lc(v,\phi,\lambda_1,\lambda_2,\lambda_3,\lambda_4)
        =&\ \langle U_q\vert \bar\De_h P_v q_v\rangle+\langle U_\phi\vert \bar\De_h P_v\phi_v\rangle\notag\\
        &-\langle W_\phi\vert \bar\De_h P_v q_v\rangle-\langle \lambda_1\vert \bar\De_h P_v q_v\rangle-\langle \lambda_2\vert \bar\De_h P_v\phi_v\rangle,
    \end{align}
    and regrouping the terms yields the first formula. The second formula follows from \cref{lemma: derivative of pushforward operator}.
\end{proof}

\subsection{Conclusion}

We can attack this problem using a primal-dual method, updates in $\lambda_i$, $i\in[5]$ are straightforward, while updating $v$ requires solving three Poisson equations, which we might be able to speed up by trying to learn the ergodic operator $(\Id-P^*_v)^{-1}$ as a function of $v$ (or equivalently the solution to the Poisson equation as a function of $(f,v)$).

The updates depend on $\nu$ only for the $\lambda_3$ update, unless we decide to collapse $q_v=\nu$ into other constraints but I don't think that will be a good idea. This means we would take $\lambda_3$ to belong to a regular function class whose regularity informs the expected regularity of the target distribution $\nu$, then we take a bunch of SGD steps on $\lambda_3$ to learn the target distribution. We take some gradient steps in $\lambda_i$, $i\in\neq 3$, compute the solutions to the Poisson equations, and then take a gradient step (or some stochastic/approximate gradient steps) in $v$ to minimise the Lagrangian.

The problem with this is that each step is probably going to be very expensive due to the integral equation resolution. On the upside, we can use a GPU friendly solver or some clever heuristics based on transport equations like in the previous work. One can hope that the convergence of the primal-dual method is very fast in terms of training, akin to a sort of second-order method, since the Poisson equations should be some kind of gradient equations for the ergodic operator $(\Id-P^*_v)^{-1}$, which is the inverse of the Hessian of the Lagrangian in $v$. It would also be possible to take SGD steps for optimising $v$ by sending individual sample points of $\mu_0$ into $q_v$ by multi-step transport (or even single step?)

%\bibliographystyle{plainnat}
\bibliography{../ngbib}
\end{document}
